\documentclass[11pt,a4paper]{article}

% ---------- Idioma y codificación ----------
\usepackage[spanish,es-tabla]{babel}
\usepackage[utf8]{inputenc}
\usepackage[T1]{fontenc}
\usepackage{lmodern}

% ---------- Formato de página ----------
\usepackage[margin=2.5cm]{geometry}
\usepackage{setspace}
\setstretch{1.15}
\usepackage{parskip}

% ---------- Tablas y listas ----------
\usepackage{booktabs}
\usepackage{tabularx}
\usepackage{array}
\usepackage{enumitem}
\newcolumntype{Y}{>{\raggedright\arraybackslash}X}
\newcolumntype{P}[1]{>{\raggedright\arraybackslash}p{#1}}

% ---------- Títulos en negrita (sin color) ----------
\usepackage{titlesec}
\titleformat{\section}
  {\Large\bfseries}
  {\thesection}{0.8em}{}
\titleformat{\subsection}
  {\large\bfseries}
  {\thesubsection}{0.6em}{}
\titlespacing*{\section}{0pt}{18pt}{8pt}

% ---------- Encabezado y pie ----------
\usepackage{fancyhdr}
\pagestyle{fancy}
\fancyhf{}
\fancyhead[L]{\small\textbf{Plan de negocios inicial}}
\fancyhead[R]{\small\textbf{CrediStep}}
\fancyfoot[C]{\small\thepage}
\renewcommand{\headrulewidth}{0.4pt}

% ---------- Caja de resumen (solo blanco y negro) ----------
\usepackage[most]{tcolorbox}
\newtcolorbox{resumen}{
  colback=white, colframe=black,
  boxrule=0.8pt, arc=0pt,
  left=8pt, right=8pt, top=6pt, bottom=6pt,
  fonttitle=\bfseries, title={La idea en pocas palabras}
}

\begin{document}
\sloppy

% ---------- Portada ----------
\begin{center}
  {\Huge\bfseries Plan de negocios inicial} \\[0.4cm]
  {\LARGE\bfseries CrediStep} \\[0.3cm]
  {\large Construye hoy el historial que necesitarás mañana} \\[0.8cm]
\end{center}

\begin{tabularx}{\textwidth}{>{\bfseries}l X}
  \toprule
  Proyecto     & CrediStep, un GPS financiero para jóvenes \\
  Asignatura   & Gestión de Innovación en TI, Pontificia Universidad Javeriana \\
  Equipo       & Beiker Romario Jaramillo Giraldo, Laura Sofía Ortiz Gómez, Sarahy Michelle Silva Pinto, Valentina Carrillo Peñuela, Miguel Ángel Gómez López \\
  Versión      & 1.0 \\
  Fecha        & Noviembre de 2026 \\
  \bottomrule
\end{tabularx}

% ---------- 1 ----------
\section{Resumen ejecutivo}

\begin{resumen}
  CrediStep es una aplicación que acompaña a los jóvenes colombianos antes de su primer
  crédito. Les hace un diagnóstico, les traza una ruta de pasos, les enseña con lecciones
  cortas y simuladores, y les muestra qué tan preparados están. No les dice qué comprar;
  les enseña cómo avanzar.
\end{resumen}

En Colombia, solo el 19,4 \% de los jóvenes de 18 a 25 años tiene un crédito formal [1]. Nuestra encuesta a 126 jóvenes mostró que el 54 \% se sentiría inseguro al pedir su primer crédito sin orientación y que el 85,7 \% quiere una herramienta que se lo explique de forma sencilla [7].

El negocio tiene dos fuentes de ingreso: una suscripción individual, CrediStep Plus, de \$20.000 al mes, y licencias para universidades. Estimamos un mercado atendible de 4,2 millones de jóvenes y una meta de 439 suscriptores de pago en el primer año, que equivalen a unos \$105 millones. Con esa meta el primer año queda prácticamente en equilibrio, y el resultado es positivo desde el segundo [8][9].

Hoy tenemos el problema validado [7], el modelo de negocio definido [12] y un MVP web en funcionamiento [11]. El siguiente paso es un piloto con la Pontificia Universidad Javeriana.

\begin{tabularx}{\textwidth}{X r}
  \toprule
  \textbf{Indicador} & \textbf{Valor} \\
  \midrule
  Mercado total (TAM), jóvenes de 18 a 30 años & 11,6 millones \\
  Mercado atendible (SAM) & 4,2 millones \\
  Suscriptores de pago en el año 1 (SOM) & 439 \\
  Ingresos del año 1 & \$105,4 millones \\
  Inversión inicial & \$5 millones \\
  Punto de equilibrio del año 1 & 440 suscriptores \\
  \bottomrule
\end{tabularx}

% ---------- 2 ----------
\section{Problema y oportunidad}

Entre los 18 y los 25 años llegan los primeros ingresos y las primeras decisiones sobre tarjetas y créditos. Casi nadie enseña a tomarlas: se aprende de la familia, del banco que vende el producto o del error. Y un error temprano, como una mora, puede cerrar puertas años después.

El crédito llega temprano, pero la orientación llega tarde. Mientras el 19,4 \% de los jóvenes de 18 a 25 años tiene crédito formal, en el grupo de 26 a 40 años la cifra sube a 37,5 \% [1][2]. Ese salto ocurre sin acompañamiento.

La oportunidad está en el tiempo que hay antes del primer crédito grande. En ese periodo se pueden formar hábitos y construir historial, y hoy ningún producto acompaña al joven paso a paso durante ese tramo.

% ---------- 3 ----------
\section{Solución}

CrediStep funciona como un GPS financiero, en cuatro pasos:

\begin{tabularx}{\textwidth}{c P{2.8cm} Y}
  \toprule
  \textbf{Paso} & \textbf{Nombre} & \textbf{Qué hace} \\
  \midrule
  1 & Diagnóstico & Preguntas cortas sobre la situación del joven. El resultado es un índice de preparación crediticia. \\
  2 & Ruta & Una lista de pasos concretos, en orden, según su situación: universitario, primer empleo o independiente. \\
  3 & Aprende y simula & Lecciones cortas con quiz, y simuladores de cuota, de uso del cupo y de atrasos. \\
  4 & Progreso & Metas de ahorro, recordatorios de pago, puntos y racha, y el avance del índice. \\
  \bottomrule
\end{tabularx}

La aplicación nunca recomienda bancos, marcas ni productos específicos. Explica tipos de productos y qué revisar antes de tomarlos. Esa neutralidad es parte de la propuesta de valor.

\subsection{Estado del producto}

El MVP es una aplicación web que ya funciona [11]. Incluye diagnóstico, ruta, lecciones, simuladores, metas, recordatorios, un asistente de preguntas frecuentes y un panel para la universidad con métricas anónimas de sus estudiantes. Sus límites actuales son tres: corre en un computador local, el plan Plus se activa en modo demostración, sin pasarela de pagos, y todavía no se ha probado con usuarios reales.

% ---------- 4 ----------
\newpage
\section{Validación}

Aplicamos una encuesta a 126 jóvenes para comprobar si el problema era real [7]. Los resultados principales fueron:

\begin{tabularx}{\textwidth}{X r}
  \toprule
  \textbf{Hallazgo} & \textbf{Resultado} \\
  \midrule
  Se sentiría inseguro al tomar su primer crédito sin orientación & 54 \% \\
  Ya evitó pedir un crédito por no saber cómo funcionaba & 35,7 \% \\
  Se ubicó en los niveles más bajos de preparación & 40,5 \% \\
  Tiene interés alto o muy alto en una herramienta de orientación & 85,7 \% \\
  \bottomrule
\end{tabularx}

La encuesta tiene dos límites que conviene decir: la mayoría de las respuestas son de jóvenes de 18 a 20 años, y no se preguntó por ocupación ni por ciudad [7]. Además, el interés declarado no equivale a disposición de pago. Por eso el plan no usa el 85,7 \% para proyectar ingresos, sino una tasa de conversión de la industria.

% ---------- 5 ----------
\section{Mercado}

El público objetivo son jóvenes de 18 a 30 años que están empezando su vida crediticia: universitarios, jóvenes con su primer empleo, trabajadores independientes y jóvenes profesionales.

\begin{tabularx}{\textwidth}{l X r}
  \toprule
  \textbf{Nivel} & \textbf{Definición} & \textbf{Usuarios} \\
  \midrule
  TAM & Jóvenes de 18 a 30 años en Colombia [3] & 11.580.726 \\
  SAM & Viven en ciudades [15], tienen internet [13][14], tienen ingresos o estudian [16], y aún no tienen crédito formal [1][2] & 4.182.110 \\
  SOM & Suscriptores de pago en el primer año & 439 \\
  \bottomrule
\end{tabularx}

El SOM se calcula así: llegar al 0,5 \% del mercado atendible da unos 20.900 registros en la versión gratuita, y de ellos el 2,1 \% pasa a pago, que es la conversión mediana de las aplicaciones con modelo freemium [4]. El alcance de 0,5 \% es un supuesto del equipo, fijado para llegar a la meta de ingresos; es la cifra que más pesa en el plan y la primera que debe medirse en el piloto. El detalle de cada filtro y de los escenarios está en el análisis de mercado, versión 2 [8].

% ---------- 6 ----------
\newpage
\section{Competencia y diferenciación}

No encontramos en Colombia un producto que acompañe al joven antes del primer crédito. Lo que existe hoy cubre partes del problema:

\begin{tabularx}{\textwidth}{P{2.9cm} Y Y}
  \toprule
  \textbf{Alternativa} & \textbf{Qué ofrece} & \textbf{Qué deja sin cubrir} \\
  \midrule
  Midatacrédito (DataCrédito Experian) & Consulta del historial y del puntaje, alertas y contenido educativo. Sus planes de pago empezaban en \$19.400 al mes en 2024 [5][6]. & Sirve sobre todo a quien ya tiene historial. No traza una ruta para quien empieza. \\
  Bancos y \textit{fintech} & Blogs y contenidos de educación financiera, gratuitos. & Quien educa es quien vende el producto. No hay seguimiento del progreso. \\
  Familia y redes sociales & Consejos cercanos y gratuitos. & Información dispersa y sin verificar. En la encuesta, familiares y amigos fueron la primera fuente elegida (23,8 \%), seguidos de los bancos (23 \%) [7]. \\
  \bottomrule
\end{tabularx}

CrediStep se diferencia en tres cosas: se enfoca en el momento anterior al primer crédito, es neutral porque no vende productos financieros, y convierte la educación en una ruta con progreso medible. El riesgo es que un competidor con más recursos, como una central de riesgo o un banco digital, lance algo parecido. Nuestra defensa es la alianza con universidades y un contenido pensado para jóvenes.

% ---------- 7 ----------
\section{Modelo de negocio}

\begin{tabularx}{\textwidth}{P{2.6cm} P{3.2cm} Y}
  \toprule
  \textbf{Plan} & \textbf{Precio} & \textbf{Qué incluye} \\
  \midrule
  Gratis & \$0 & Diagnóstico, ruta, lecciones esenciales, simuladores básicos y hasta dos metas. Es la puerta de entrada. \\
  CrediStep Plus & \$20.000 al mes o \$200.000 al año & Simulador de puntaje, tabla de amortización completa, lecciones avanzadas y metas ilimitadas. \\
  Licencia institucional & Por definir & Plus para los estudiantes de la universidad y un panel con métricas anónimas. \\
  \bottomrule
\end{tabularx}

El cliente de la suscripción es el mismo joven. El cliente de la licencia es la universidad, que la ofrece a sus estudiantes como parte de su bienestar financiero. La licencia todavía no tiene precio, así que este plan no cuenta sus ingresos: todo lo que aporte es adicional a lo proyectado. Los precios y la regla de la licencia (el código institucional da Plus sin costo) vienen de la aplicación [11] y del análisis de mercado [8]; el modelo completo está en el lienzo Canvas [12].

Los recursos clave son el equipo de desarrollo, el contenido educativo y la alianza con universidades. Los costos principales son la nómina, la infraestructura en la nube y el mercadeo.

% ---------- 8 ----------
\section{Estrategia de mercadeo y ventas}

\begin{itemize}[leftmargin=1.5em]
  \item \textbf{Canal del producto.} La aplicación web, que ya se adapta al celular. Después, las tiendas de aplicaciones de Android y iOS.
  \item \textbf{Universidades.} Empezamos con la Javeriana. La alianza da la credibilidad que una aplicación nueva no tiene por sí sola, y concentra a muchos jóvenes en un solo lugar.
  \item \textbf{Redes sociales y creadores de contenido.} Contenido corto y cercano, en los formatos que los jóvenes ya usan.
  \item \textbf{Conversión a pago.} El plan gratuito resuelve lo esencial. Plus se ofrece cuando el usuario ya avanzó en su ruta y quiere simular su puntaje o planear un crédito.
  \item \textbf{Retención.} Recordatorios de pago, metas, puntos y racha, para que el joven vuelva cada semana.
\end{itemize}

El presupuesto de mercadeo es de \$5 millones al año en los dos primeros años y de \$10 millones desde el tercero. Es un presupuesto pequeño frente a la meta de 20.900 registros, así que la adquisición depende sobre todo de canales sin costo directo: campus, voz a voz y contenido propio.

% ---------- 9 ----------
\section{Equipo y operación}

\begin{tabularx}{\textwidth}{l X}
  \toprule
  \textbf{Integrante} & \textbf{Rol} \\
  \midrule
  Beiker Jaramillo & Scrum Master \\
  Laura Ortiz & Desarrolladora \textit{full-stack} \\
  Sarahy Silva & Diseñadora UX/UI \\
  Valentina Carrillo & Ingeniera de QA \\
  Miguel Gómez & Investigación y validación \\
  \bottomrule
\end{tabularx}

Somos cinco estudiantes de Ingeniería de Sistemas. Conocemos el problema porque lo vivimos, y tenemos la capacidad técnica para construir el producto. Nos falta una persona con formación financiera que revise el contenido sobre tasas y normas antes de publicarlo; la buscaremos como asesora durante el piloto.

La aplicación está construida con FastAPI en el servidor y Vue en la interfaz. Para operar con usuarios reales debe pasar a la nube, con una base de datos PostgreSQL, y sumar una pasarela de pagos.

% ---------- 10 ----------
\section{Plan financiero}

\subsection{Supuestos}

\begin{tabularx}{\textwidth}{X r}
  \toprule
  \textbf{Supuesto} & \textbf{Valor} \\
  \midrule
  Ingreso anual por suscriptor Plus (\$20.000 al mes) & \$240.000 \\
  Suscriptores de pago en el año 1 & 439 \\
  Crecimiento de suscriptores en los años 2, 3 y 4 & 10 \%, 30 \% y 60 \% \\
  Inversión inicial, aportada por el equipo & \$5 millones \\
  Nómina del año 1 & \$90 millones \\
  Infraestructura y tiendas de aplicaciones, año 1 & \$10,6 millones \\
  Mercadeo, año 1 & \$5 millones \\
  \bottomrule
\end{tabularx}

Los costos y las tasas de crecimiento vienen de la hoja de simulación financiera, versión 2 [9]. Los suscriptores y el precio vienen del análisis de mercado [8]. La explicación completa y el análisis de sensibilidad están en el contexto financiero [10].

\subsection{Proyección a cuatro años}

\begin{tabularx}{\textwidth}{X r r r r}
  \toprule
  \textbf{Concepto (millones de COP)} & \textbf{Año 1} & \textbf{Año 2} & \textbf{Año 3} & \textbf{Año 4} \\
  \midrule
  Suscriptores de pago & 439 & 483 & 628 & 1.004 \\
  Ingresos por suscripciones & 105,4 & 115,9 & 150,7 & 241,0 \\
  Nómina & 90,0 & 95,0 & 98,8 & 102,5 \\
  Infraestructura y tiendas & 10,6 & 11,2 & 11,6 & 12,1 \\
  Mercadeo & 5,0 & 5,0 & 10,0 & 10,0 \\
  Costos y gastos totales & 105,6 & 111,1 & 120,5 & 124,6 \\
  \textbf{Resultado antes de depreciación e impuestos} & \textbf{$-$0,2} & \textbf{4,8} & \textbf{30,2} & \textbf{116,4} \\
  \bottomrule
\end{tabularx}

El punto de equilibrio del primer año es de 440 suscriptores: \$105,6 millones de costos divididos entre \$240.000 por suscriptor. La meta de 439 queda justo en ese borde, sin margen. El resultado acumulado de los cuatro años es de unos \$151 millones.

\subsection{Lo que esta proyección no incluye}

\begin{itemize}[leftmargin=1.5em]
  \item Los ingresos por licencias a universidades, porque no tienen precio definido.
  \item Las comisiones de la pasarela de pagos y de las tiendas de aplicaciones. La depreciación y el impuesto de renta solo se calculan en la hoja financiera [9], no en esta tabla.
  \item La cancelación de suscripciones. Suponemos que cada suscriptor paga los doce meses; si paga menos, el ingreso baja en proporción.
  \item La hoja de simulación versión 2 ya usa estos supuestos [9], pero cuenta cada suscriptor con doce meses de pago. Si los suscriptores entran poco a poco, el recaudo de los primeros doce meses sería de unos \$45,4 millones [9][10].
\end{itemize}

% ---------- 11 ----------
\section{Riesgos}

\begin{tabularx}{\textwidth}{Y Y}
  \toprule
  \textbf{Riesgo} & \textbf{Cómo lo manejamos} \\
  \midrule
  No llegar a 20.900 registros con un presupuesto de mercadeo pequeño & Empezar por la universidad y medir en el piloto el costo real de cada registro. \\
  Que menos del 2,1 \% pase a pago & Probar precios y el momento de la oferta de Plus; apoyarse en las licencias. \\
  Contenido financiero con errores & Revisión de una persona con formación financiera antes de publicar. \\
  Manejo de datos personales & Pedir solo los datos necesarios y mostrar a la universidad métricas anónimas. \\
  Un competidor con más recursos & Alianzas con universidades y contenido pensado para jóvenes. \\
  Que el equipo no pueda dedicarse de tiempo completo & Avanzar por etapas y condicionar la nómina a los resultados del piloto. \\
  \bottomrule
\end{tabularx}

% ---------- 12 ----------
\section{Hoja de ruta}

\begin{tabularx}{\textwidth}{P{4.2cm} Y}
  \toprule
  \textbf{Etapa} & \textbf{Qué haremos} \\
  \midrule
  Noviembre de 2026 & Pitch final y entrega del MVP documentado. \\
  Primer semestre de 2027 & Piloto con la Javeriana, pruebas con usuarios reales y definición del precio de la licencia. \\
  Segundo semestre de 2027 & Pasarela de pagos, despliegue en la nube y notificaciones de recordatorios. \\
  2028 & Versión móvil en tiendas de aplicaciones y nuevas universidades. \\
  \bottomrule
\end{tabularx}

El piloto debe responder tres preguntas: cuántos jóvenes se registran, cuántos siguen usando la aplicación después del primer mes y cuántos pagarían por Plus. Con esas tres cifras se ajusta este plan.

% ---------- 13 ----------
\section{Referencias}

\begin{enumerate}[label={[\arabic*]},leftmargin=2em]
  \item Superintendencia Financiera de Colombia y Banca de las Oportunidades. Reporte de Inclusión Financiera 2025.
  \item Banca de las Oportunidades y Superintendencia Financiera de Colombia. Reporte de Inclusión Financiera 2024.
  \item DANE. Proyecciones de población 2026, consultadas a través de telencuestas.com, que reproduce las cifras del DANE.
  \item RevenueCat. State of Subscription Apps 2026.
  \item Experian. Midatacrédito, aplicación de acompañamiento financiero. Sitio corporativo de Experian (experianplc.com), sección de salud financiera.
  \item Infobae (12 de octubre de 2024). Más de 600.000 colombianos revisan mensualmente su historial crediticio.
  \item CrediStep. Artefactos de Validación: encuesta a 126 jóvenes (\textit{Artefacto\_De\_Validacion.pdf}), Sprint 3.
  \item CrediStep. Análisis de mercado: TAM, SAM y SOM, versión 2.0 (\textit{Definicion\_de\_mercado\_V2.pdf}), Sprint 5. Noviembre de 2026.
  \item CrediStep. Proyección financiera, versión 2 (\textit{Proyeccion\_Financiera\_V2.xlsx}), Sprint 7.
  \item CrediStep. Contexto financiero (\textit{Contexto\_Financiero.pdf}), Sprint 7.
  \item CrediStep. Cómo funciona la aplicación (\textit{CrediStep\_Como\_Funciona\_la\_Aplicacion.pdf}), Sprint 9, y README de \textit{CrediStep-App}.
  \item CrediStep. Lienzo Canvas (\textit{Lienzo\_Canvas.pdf}), Sprint 4.
  \item DANE. Encuesta de Calidad de Vida 2025: 82,3 \% de las personas de 5 años o más usa internet. Citado en Forbes Colombia.
  \item MinTIC y DANE. ENTIC: 97,5 \% de los usuarios de internet se conecta desde el celular.
  \item DANE. Censo Nacional de Población y Vivienda 2018: 77 \% de la población vive en cabeceras municipales. Citado en La República.
  \item DANE. Gran Encuesta Integrada de Hogares, mercado laboral de la juventud, mayo a julio de 2025: 20,6 \% de jóvenes que ni estudian ni trabajan. Citado en La República.
\end{enumerate}

Las referencias [7] a [12] son documentos del proyecto y están en el repositorio: Artefacto de Validación (Sprint 3), lienzo Canvas (Sprint 4), análisis de mercado versión 2 (Sprint 5), hoja y contexto financiero (Sprint 7) y documentación de la aplicación (Sprint 9).

\end{document}
